% Пояснювальна записка до курсової роботи. Збирання: xelatex main.tex (двічі).
\documentclass[14pt,a4paper]{extarticle}

\usepackage{fontspec}
\usepackage{polyglossia}
\setdefaultlanguage{ukrainian}
\setotherlanguage{english}
\setmainfont{FreeSerif}
\setsansfont{DejaVu Sans}
\setmonofont[Scale=0.82]{DejaVu Sans Mono}
\newfontfamily\cyrillicfont{FreeSerif}
\newfontfamily\cyrillicfonttt[Scale=0.82]{DejaVu Sans Mono}

\usepackage[left=30mm,right=15mm,top=20mm,bottom=20mm,headsep=6mm,headheight=14pt]{geometry}
\usepackage{setspace}
\onehalfspacing
\setlength{\parindent}{12.5mm}
\usepackage{indentfirst}
\usepackage{graphicx}
\usepackage{float}
\usepackage{amsmath}
\usepackage{booktabs,longtable,array,multirow,tabularx}
\usepackage[labelsep=endash,font=normal,justification=centering]{caption}
\usepackage{enumitem}
\usepackage{fancyhdr}
\usepackage{titlesec}
\usepackage{tocloft}
\usepackage{totcount}
\usepackage{lastpage}
\usepackage{xurl}
\usepackage[hidelinks]{hyperref}
\usepackage{listings}
\usepackage{etoolbox}

% ---------- колонтитули: номер сторінки праворуч угорі (ДСТУ 3008:2015)
\pagestyle{fancy}
\fancyhf{}
\fancyhead[R]{\thepage}
\renewcommand{\headrulewidth}{0pt}
\fancypagestyle{plain}{\fancyhf{}\fancyhead[R]{\thepage}\renewcommand{\headrulewidth}{0pt}}

% ---------- заголовки
\titleformat{\section}{\normalfont\bfseries\centering}{\thesection}{1ex}{\MakeUppercase}
\titleformat{\subsection}{\normalfont\bfseries}{\thesubsection}{1ex}{}
\titleformat{\subsubsection}{\normalfont\bfseries}{\thesubsubsection}{1ex}{}
\titlespacing*{\section}{0pt}{0pt}{18pt}
\titlespacing*{\subsection}{\parindent}{14pt}{8pt}
\titlespacing*{\subsubsection}{\parindent}{10pt}{6pt}
\newcommand{\sectionbreak}{\clearpage}
% структурні елементи без номера
\newcommand{\unsection}[1]{\section*{#1}\addcontentsline{toc}{section}{#1}}

% ---------- зміст
\gappto\captionsukrainian{\renewcommand{\contentsname}{ЗМІСТ}}
\renewcommand{\cfttoctitlefont}{\hfill\bfseries}
\renewcommand{\cftaftertoctitle}{\hfill}
\renewcommand{\cftsecfont}{\normalfont}
\renewcommand{\cftsecpagefont}{\normalfont}
\renewcommand{\cftsecleader}{\cftdotfill{\cftdotsep}}
\setlength{\cftbeforesecskip}{2pt}
\setcounter{tocdepth}{2}

% ---------- нумерація рисунків, таблиць і формул у межах розділу
\counterwithin{figure}{section}
\counterwithin{table}{section}
\counterwithin{equation}{section}
\captionsetup[figure]{name=Рисунок,skip=6pt}
\captionsetup[table]{name=Таблиця,justification=raggedright,singlelinecheck=false,skip=4pt}

% лічильники для реферату
\regtotcounter{figure}
\newtotcounter{figs}
\newtotcounter{tabs}
\newtotcounter{refs}
\pretocmd{\figure}{\stepcounter{figs}}{}{}
\pretocmd{\table}{\stepcounter{tabs}}{}{}
\pretocmd{\longtable}{\stepcounter{tabs}}{}{}

\newcommand{\fig}[3][0.95\textwidth]{%
  \begin{figure}[H]\centering
  \includegraphics[width=#1,height=0.80\textheight,keepaspectratio]{figures/#2}
  \caption{#3}\label{fig:#2}\end{figure}}

\setlist{nosep,leftmargin=*,labelindent=\parindent,itemindent=0pt}
\setlist[itemize]{label=--}
\lstset{basicstyle=\ttfamily\small,columns=fullflexible,breaklines=true,
  frame=single,xleftmargin=2mm,aboveskip=6pt,belowskip=6pt}
\newcommand{\code}[1]{\texttt{#1}}
\renewcommand{\arraystretch}{1.15}
\sloppy
\makeatletter\newcommand{\tabinput}[1]{\@@input #1 }\makeatother

\begin{document}

% ======================================================================== ТИТУЛЬНИЙ АРКУШ
\begin{titlepage}
\begin{singlespace}
\centering
Міністерство освіти і науки України\\
Національний технічний університет України\\
«Київський політехнічний інститут імені Ігоря Сікорського»\\
Навчально-науковий інститут прикладного системного аналізу\\
Кафедра системного проєктування

\vspace{55mm}
{\bfseries КУРСОВА РОБОТА}\\[4mm]
з дисципліни «Паралельні обчислення»\\[2mm]
на тему:\\[2mm]
{\bfseries «Дослідження та реалізація паралельних обчислень\\
на прикладі побудови пошукової системи з інтеграцією LLM»}

\vspace{40mm}
\begin{flushright}
\begin{minipage}{0.55\textwidth}
Виконав: студент групи КН-32\\
Майстренко Ілля Ігорович\\[6mm]
Керівник: доктор філософії,\\ ст. викладач кафедри СП\\
Яременко Вадим Сергійович
\end{minipage}
\end{flushright}

\vfill
Київ – 2026
\end{singlespace}
\end{titlepage}
\setcounter{page}{2}

\section*{РЕФЕРАТ}

Пояснювальна записка до курсової роботи: \pageref{LastPage}~с., \total{figs}~рис.,
\total{tabs}~табл., 3~додатки, \total{refs}~джерело.

\textit{Об'єкт дослідження} -- паралельна побудова інвертованого індексу та паралельна обробка
пошукових запитів у пошуковій системі з розширенням запитів великою мовною моделлю.

\textit{Мета роботи} -- розробити контейнеризовану пошукову систему мовою Java без сторонніх
бібліотек (власний пул потоків, власна потокобезпечна структура індексу, власний HTTP/1.1- і
WebSocket-сервер, розширення запитів локальною моделлю qwen2.5:1.5b через Ollama) та
експериментально дослідити її масштабованість.

\textit{Методи дослідження}: вимірювання часу побудови індексу на 1--16 потоках для двох
реалізацій спільної структури (неблокуючої з хеш-партиціями і блокувальної з lock striping) і
для 1, 8, 64 та 256 партицій, по три прогони з випадковим чергуванням конфігурацій; аналіз
прискорення за законом Амдала та метрикою Карпа -- Флатта; розклад часу за фазами;
порівняння двох версій токенізації; навантажувальне тестування пошуку для 1--64 одночасних
клієнтів; вимірювання пам'яті під час повторних перебудов; модульне та наскрізне тестування.

\textit{Результати}. На корпусі з 16\,793 документів (1,62~ГБ) медіанний час побудови індексу
скоротився з 10,12~с на одному потоці до 2,09~с на 10 потоках (прискорення 4,85) на
процесорі з 6 ядрами і 12 апаратними потоками. Заміна токенізації, що створювала рядок на
кожне слово, байтовою хеш-таблицею задачі скоротила час у 1,28--1,48 раза за будь-якої
кількості потоків. З однією партицією блокувальний варіант на 12 потоках сумарно очікує на
блокування 22,4~с, а неблокуючий не очікує, але обидва впираються в послідовне злиття
(3,4--3,5~с); від 8 партицій синхронізація перестає бути вузьким місцем. Пропускна здатність
пошуку сягає 12\,004 запитів/с на 16 клієнтах; з розширенням через LLM вона обмежена моделлю
на рівні 2,3--2,4 запиту/с, а кешування розширень повертає її до 8\,634 запитів/с.

\bigskip
\noindent ПАРАЛЕЛЬНІ ОБЧИСЛЕННЯ, ПУЛ ПОТОКІВ, ІНВЕРТОВАНИЙ ІНДЕКС, НЕБЛОКУЮЧА СИНХРОНІЗАЦІЯ,
CAS, ЗАКОН АМДАЛА, МЕТРИКА КАРПА -- ФЛАТТА, BM25, WEBSOCKET, LLM, OLLAMA, DOCKER

\clearpage

\tableofcontents
\unsection{ПЕРЕЛІК СКОРОЧЕНЬ}

\begin{description}[leftmargin=0pt,labelindent=\parindent,labelsep=0.5em,itemsep=0pt,
  font=\normalfont]
\item[API] -- прикладний програмний інтерфейс (Application Programming Interface);
\item[BM25] -- функція ранжування Okapi Best Matching 25;
\item[CAS] -- атомарна операція «порівняти й замінити» (Compare-And-Swap);
\item[CLI] -- інтерфейс командного рядка (Command-Line Interface);
\item[CPU] -- центральний процесор (Central Processing Unit);
\item[GC] -- збирач сміття віртуальної машини Java (Garbage Collector);
\item[HTTP] -- протокол передавання гіпертексту (HyperText Transfer Protocol);
\item[IDF] -- обернена документна частота (Inverse Document Frequency);
\item[JDK, JRE] -- комплект розробника / середовище виконання Java;
\item[JIT] -- компіляція під час виконання (Just-In-Time);
\item[JMM] -- модель пам'яті Java (Java Memory Model);
\item[JSON] -- текстовий формат обміну даними (JavaScript Object Notation);
\item[LLM] -- велика мовна модель (Large Language Model);
\item[LRU] -- витіснення найдавніше використаного (Least Recently Used);
\item[REST] -- архітектурний стиль вебінтерфейсів (Representational State Transfer);
\item[RFC] -- серія документів-специфікацій інтернету (Request for Comments);
\item[RSS] -- резидентна пам'ять процесу (Resident Set Size);
\item[SMT] -- одночасна багатопотоковість ядра процесора (Simultaneous Multithreading);
\item[TCP] -- протокол керування передаванням (Transmission Control Protocol);
\item[TF] -- частота терма в документі (Term Frequency);
\item[UI] -- інтерфейс користувача (User Interface);
\item[UML] -- уніфікована мова моделювання (Unified Modeling Language);
\item[WS] -- протокол WebSocket (RFC 6455);
\item[WSL2] -- підсистема Windows для Linux, версія 2;
\item[ОЗП] -- оперативний запам'ятовувальний пристрій.
\end{description}

\unsection{ВСТУП}

\subsection*{Актуальність, мета та постановка задачі}

Повнотекстовий пошук є типовою задачею, де паралелізм потрібен на двох рівнях одночасно.
По-перше, побудова інвертованого індексу над великою колекцією -- це обчислювально важка
пакетна робота: кожен байт тексту треба прочитати, розбити на слова й врахувати у
спільній структурі даних. По-друге, готовий індекс обслуговує багатьох клієнтів одночасно, і
сервер має обробляти їхні запити паралельно, не блокуючи один одного. Основне завдання
паралельної обробки -- скорочення часу виконання програми [1]; у розробленій системі
побудова індексу корпусу з 16\,793 документів на одному потоці триває 10,1~с, а на десяти --
2,1~с (підрозділ~3.4). Третій чинник -- інтеграція великої мовної моделі для розширення
запитів: модель на CPU відповідає за сотні мілісекунд і стає послідовною ділянкою конвеєра,
яку доводиться ізолювати від решти системи.

\textit{Метою роботи} є дослідження та практична реалізація паралельних обчислень на
прикладі пошукової системи з інтеграцією LLM: паралельної побудови інвертованого індексу за
допомогою власного пулу потоків, потокобезпечного доступу до спільної структури індексу та
одночасної обробки запитів від багатьох клієнтів.

\textit{Об'єкт дослідження} -- процеси паралельної побудови інвертованого індексу та
паралельної обробки пошукових запитів. \textit{Предмет дослідження} -- механізми
синхронізації потоків у пошуковій системі (блокувальні та неблокуючі способи злиття
результатів у спільну структуру, безблокувальне читання незмінного знімка індексу, обмеження
паралельних звернень до моделі) та їхній вплив на час побудови, пропускну здатність і
використання пам'яті. \textit{Методи дослідження} -- повторні вимірювання з випадковим
чергуванням конфігурацій, аналіз прискорення за законом Амдала та метрикою Карпа -- Флатта
(підрозділ~3.1), розклад часу за фазами, порівняльний експеримент «до/після» оптимізації,
модульне, стресове та наскрізне тестування.

Відповідно до завдання система має задовольняти такі вимоги:
\begin{itemize}
\item бекенд компільованою мовою (обрано Java 21), два різні клієнти -- консольний і веб;
\item запуск однією командою \code{docker compose up}: сервіс пошукового сервера і сервіс
Ollama, ендпоінт \code{GET /api/health};
\item самостійно реалізовані пул потоків і структура інвертованого індексу з паралельним
доступом, кількість потоків побудови задається параметром \code{threads};
\item фіксований REST-контракт: \code{GET /api/health}, \code{POST /api/index},
\code{POST /api/search} (режим \code{keyword}), \code{GET /api/stats}, формат JSON;
\item розширення запиту ключовими словами та синонімами через Ollama \code{/api/generate};
\item WebSocket \code{/ws/search}, що передає етапи \code{query\_expansion},
\code{keyword\_search}, \code{done} окремими повідомленнями з полями \code{stage} і
\code{data};
\item текстовий датасет англійською мовою обсягом не менше 1~ГБ і не менше 10\,000 документів;
\item графіки часу побудови індексу від кількості потоків і пропускної здатності від
кількості клієнтів, аналіз вузьких місць, UML-діаграми, інструкція зі збирання, історія
комітів у стилі Conventional Commits.
\end{itemize}

\subsection*{Обрані технології}

\textit{Java 21 без сторонніх бібліотек.} Мова надає модель пам'яті з чітко визначеним
відношенням happens-before [15], атомарні змінні з операцією CAS (\code{AtomicReference},
\code{AtomicBoolean}) та примітиви \code{ReentrantLock}/\code{Condition} [4] -- саме ті засоби,
порівняння яких є предметом роботи. Усі компоненти, що зазвичай беруть із бібліотек, --
пул потоків, інвертований індекс, HTTP/1.1-сервер, протокол WebSocket, розбір JSON,
навантажувальний клієнт і навіть тестовий фреймворк -- реалізовано в проєкті. Це
спрощує перевірку відсутності станів гонки (весь код синхронізації видимий) і дозволяє
збирати проєкт лише засобами JDK (\code{javac} і \code{jar}) без доступу до репозиторіїв
залежностей.

\textit{Docker і Docker Compose} [17] забезпечують відтворюване розгортання: сервер, Ollama та
два одноразові сервіси ініціалізації (розпакування корпусу і завантаження моделі)
піднімаються однією командою.

\textit{Ollama і модель qwen2.5:1.5b} [16, 21]. Модель обрано як компроміс між якістю
синонімів і затримкою на CPU: вона дає змістовні розширення (табл.~3.10) з медіанною
затримкою 0,64~с (підрозділ~3.9); більші моделі тієї самої родини потребують пропорційно
більше обчислень на кожен згенерований токен і на CPU відповідали б повільніше.

\textit{Датасет} -- архів книжок Project Gutenberg [20] обсягом 1,70~ГБ (2\,809 файлів). Оскільки
книжок менше за потрібні 10\,000 документів, кожну книжку розрізано по межах абзаців на
документи приблизно по 100~КБ (підрозділ~2.1), що дало 16\,793 документи обсягом 1,62~ГБ.
У силабусі в одному місці вказано обсяг «не менше 2~ГБ», а у вимогах до датасету -- «не менше
1~ГБ»; робота виконує вимогу розділу «Вимоги до датасету».

\textit{Допоміжні засоби}: PlantUML для UML-діаграм, Python з matplotlib для побудови
графіків із CSV-результатів, Git і GitHub для керування версіями. Повідомлення комітів
оформлено за специфікацією Conventional Commits [18], код -- за Google Java Style [19].

\section{Модель системи}

\subsection{Діаграма прецедентів}

Систему використовують чотири категорії акторів (рис.~\ref{fig:use_case.png}): користувач
веб-інтерфейсу, користувач CLI-клієнта, дослідник, який запускає навантажувальні тести, і
оркестратор Docker, що перевіряє стан контейнера. Зовнішня система Ollama бере участь у
розширенні запиту та в перевірці стану. Прецеденти відповідають ендпоінтам фіксованого
контракту: побудова індексу (UC1), пошук через REST (UC2), пошук зі стрімінгом етапів через
WebSocket (UC3), перегляд статистики (UC4) і перевірка стану (UC5). Розширення запиту
включене в UC3 завжди (його можна вимкнути полем \code{expand:false}), а в UC2 розширює
його лише за явного прапорця \code{expand:true} -- за замовчуванням REST-пошук працює у
чистому режимі \code{keyword}.

\fig[0.98\textwidth]{use_case.png}{Діаграма прецедентів}

\subsection{Діаграма класів}

Код поділено на кореневий пакет і вісім підпакетів (рис.~\ref{fig:class_overview.png}). Пакет \code{core} містить
обчислювальне ядро: власний пул потоків \code{ThreadPool}, оркестратор побудови
\code{Indexer}, інтерфейс \code{IndexBuilder} з двома реалізаціями спільної структури
(\code{PartitionedIndexBuilder} -- неблокуюча, \code{StripedIndexBuilder} -- блокувальна),
незмінний індекс \code{InvertedIndex} з таблицею документів \code{DocumentTable},
неблокуючий стек \code{TreiberStack}, список входжень \code{PostingList}, токенізатор і
байтову таблицю термів задачі \code{BatchTermTable}. Пакет \code{net} реалізує мережевий
рівень (\code{HttpServer}, \code{HttpRequest}, \code{HttpResponse}, \code{WebSocket}), пакет
\code{api} -- маршрутизацію і спільний для REST і WebSocket конвеєр пошуку
(\code{Router}, \code{SearchService}, \code{WsSearchHandler}), пакет \code{search} --
ранжування BM25 і сніпети, пакет \code{llm} -- клієнт Ollama і розширювач запитів. Класи
\code{SearchCli} і \code{Bench} є окремими точками входу того самого jar-файла. Детальні
діаграми класів із полями синхронізації наведено в додатку~А (рис.~А.1--А.2).

\fig[0.98\textwidth]{class_overview.png}{Діаграма класів (огляд за пакетами)}

\subsection{Діаграми послідовності}

\textit{Побудова індексу} (рис.~\ref{fig:seq_index.png}). \code{Router} перевіряє параметри
\code{threads} (1--256), \code{partitions} (1--4096), \code{batch\_size} і назву реалізації;
ціле число приймається також у вигляді \code{4.0} або \code{"4"}, інші значення дають 400.
\code{Indexer} захоплює \code{buildLock} методом \code{tryLock}: паралельний запит на
перебудову не стає в чергу, а одразу отримує 409. Далі на викликаючому потоці формується
відсортований список файлів -- позиція файла в ньому є ідентифікатором документа, тому
результат не залежить від кількості потоків. Список ріжеться на батчі по 32 документи, і
кожен батч стає окремою задачею пулу. Задача читає й токенізує свої документи у приватну
таблицю без синхронізації і передає її будівнику; будівник розбиває таблицю на фрагменти за
партиціями і зливає їх неблокуючим протоколом (підрозділ~2.5). Після завершення всіх задач
по одній задачі на партицію сортує і «заморожує» списки входжень, а новий знімок індексу
атомарно підміняє попередній.

\fig{seq_index.png}{Діаграма послідовності: побудова індексу}

\textit{Пошук через REST} (рис.~\ref{fig:seq_search_rest.png}). Некоректний JSON, порожній
запит або непідтримуваний режим дають 400. За прапорця \code{expand} запит спершу
розширюється: результат береться з LRU-кешу, або з уже виконуваного однакового запиту, або
з нового виклику моделі в межах семафора і дедлайну; помилка чи тайм-аут моделі дають
порожній список термів, і пошук продовжується. Знімок індексу береться вже \textit{після}
розширення, тому запит, що чекає на модель, не утримує в пам'яті індекс, який паралельна
перебудова вже замінила. Ранжування BM25 і формування сніпетів виконуються без блокувань.

\fig{seq_search_rest.png}{Діаграма послідовності: пошук через REST API}

\textit{Пошук через WebSocket} (рис.~\ref{fig:seq_search_ws.png}). Після рукостискання сокет
передається з пулу HTTP-потоків у пул WebSocket-сеансів. В одному сеансі можна надіслати
кілька запитів поспіль; на кожен сервер відповідає трьома повідомленнями-етапами, а для
помилок окремого запиту (некоректний JSON, індекс ще не побудовано) надсилає повідомлення з
етапом \code{error}, не закриваючи з'єднання. Порушення протоколу закривають сеанс із кодом
RFC~6455.

\fig{seq_search_ws.png}{Діаграма послідовності: пошук через WebSocket}

\textit{Статистика та перевірка стану} (рис.~\ref{fig:seq_stats.png},
\ref{fig:seq_health.png}). Обидва запити лише читають знімок індексу через volatile-посилання.
Перевірка стану зондує Ollama запитом \code{GET /api/tags} з тайм-аутом 1~с і кешує
результат на 2~с, щоб часті звернення healthcheck не створювали навантаження на модель.

\fig[0.85\textwidth]{seq_stats.png}{Діаграма послідовності: статистика}
\fig[0.85\textwidth]{seq_health.png}{Діаграма послідовності: перевірка стану}

\subsection{Діаграма розгортання}

Система розгортається на одній машині (рис.~\ref{fig:deployment.png}) у проєкті Docker
Compose \code{pocw-search}. Сервіс \code{server} (образ на базі \code{eclipse-temurin:21-jre},
порт 8080) містить jar-файл, статичний Web UI та обгортки \code{search} і \code{bench}.
Сервіс \code{ollama} (порт 11434) зберігає ваги моделі в іменованому томі \code{ollama\_data},
тому повторні запуски не завантажують модель знову. Два одноразові сервіси виконуються під
час першого запуску: \code{ollama-init} завантажує модель, \code{dataset-init} розрізає архів
корпусу на документи в іменованому томі \code{corpus}, який сервер монтує тільки для
читання. Сервер стартує після завершення \code{dataset-init}, але не чекає на модель: до її
завантаження пошук працює без розширення.

\fig[0.82\textwidth]{deployment.png}{Діаграма розгортання (топологія Docker Compose)}

\section{Опис реалізації}

\subsection{Підготовка корпусу}

Скрипт \code{tools/prepare\_dataset.py} читає книжки з архіву у відсортованому порядку,
вирізає ліцензійні шапку й футер Project Gutenberg (між маркерами
\code{*** START OF ...} і \code{*** END OF ...}), нормалізує кінці рядків і пакує абзаци в
документи розміром близько 100~КБ. Абзац, довший за ліміт, ріжеться за останнім переведенням
рядка або пробілом до межі, тож слова не розриваються; короткий «хвіст» книжки (менше чверті
ліміту) приєднується до попереднього документа. Скрипт детермінований: повторний запуск дає
ті самі файли. У Docker Compose його виконує одноразовий сервіс \code{dataset-init}, записуючи
результат в іменований том \code{corpus} на файловій системі Linux.

Розміщення корпусу в томі, а не в каталозі диска Windows, -- результат першого ж вимірювання.
Коли сервер читав документи безпосередньо з \code{D:\textbackslash} через механізм спільних
каталогів Docker Desktop (міст Windows $\leftrightarrow$ WSL2), побудова на 12 потоках
тривала 42,3~с, з яких 19,3~с займав лише обхід каталогу, а сумарний час потоків на читання
файлів становив 238~с: обчислення фактично стояли в черзі до файлової системи. Після
перенесення корпусу в том обхід каталогу скоротився до 0,13~с, а вся побудова -- до 4,0~с.
Характеристики отриманого корпусу наведено в табл.~\ref{tab:corpus}.

\begin{table}[H]
\caption{Характеристики корпусу}\label{tab:corpus}
\centering\begin{tabularx}{\textwidth}{Xr}
\toprule
Показник & Значення\\
\midrule
Книжок в архіві D1.7GB.zip & 2\,809\\
Документів після розрізання & 16\,793\\
Обсяг документів & 1\,623\,765\,577 байт (1,62~ГБ)\\
Цільовий розмір документа & $\approx$100~КБ\\
Унікальних термів у словнику & 840\,954\\
Записів у списках входжень (пар «терм -- документ») & 33\,197\,116\\
Середня довжина документа після токенізації & 5\,501 терм\\
\bottomrule
\end{tabularx}
\end{table}

\subsection{Пул потоків}

Клас \code{ThreadPool} -- фіксований набір робочих потоків і необмежена FIFO-черга задач
\code{ArrayDeque<Runnable>}, захищена одним \code{ReentrantLock}. Вільні потоки чекають на
умові \code{notEmpty} у циклі з перевіркою предиката, що захищає від хибних пробуджень і
виключає активне очікування. Метод \code{submit} загортає задачу в \code{FutureTask} і
повертає \code{Future}: виняток задачі зберігається в ньому й повертається викликачу через
\code{ExecutionException}, а робочий потік продовжує роботу. Стан пулу змінюється за схемою
\code{RUNNING} $\rightarrow$ \code{SHUTDOWN} $\rightarrow$ \code{TERMINATED}: \code{shutdown()}
перестає приймати задачі, але дає виконати вже поставлені; \code{shutdownNow()} додатково
повертає невиконані задачі й перериває потоки; \code{close()} виконує коректне завершення з
приєднанням усіх потоків, що дозволяє використовувати пул у конструкції try-with-resources.

Інваріанти пулу: черга і стан читаються та змінюються лише під замком; задача виймається з
черги рівно одним потоком, тож виконується рівно один раз; \code{Future.get()} встановлює
відношення happens-before між діями задачі та кодом, який дочекався її результату [15]. Ця
остання властивість використовується в побудові індексу для безпечної передачі даних між
фазами без додаткових бар'єрів. У сервері працюють три види пулів: пул HTTP-з'єднань (64
потоки), пул WebSocket-сеансів (32 потоки) і пул побудови індексу, який створюється на кожен
запит індексації з кількістю потоків \code{threads}.

\subsection{Структура інвертованого індексу}

Інвертований індекс зіставляє кожному терму список документів, у яких він трапляється [3].
У реалізації кожен запис списку (posting) упаковано в одне 64-бітне число: старші 32 біти --
ідентифікатор документа, молодші -- частота терма в ньому. Упаковка вдвічі зменшує кількість
об'єктів порівняно з парою масивів і перетворює сортування за документом на звичайне
\code{Arrays.sort(long[])}, оскільки ідентифікатори невід'ємні. Словник поділено на $P$
хеш-партицій (за замовчуванням 64); партиція терма визначається хешуванням Фібоначчі
$h \cdot \mathtt{0x9E3779B9}$ з подальшим взяттям залишку, що рівномірно розподіляє схожі
рядки. Кожна партиція -- звичайна \code{HashMap<String, long[]>}.

Готовий \code{InvertedIndex} незмінний: після побудови жодна партиція не модифікується. Він
публікується разом зі статистикою побудови як об'єкт \code{Snapshot} через
\code{AtomicReference}. Запис у volatile-посилання після завершення всіх фаз і його читання
пошуковим потоком утворюють пару release/acquire, тому всі записи, зроблені під час побудови,
видимі читачеві без жодних блокувань [15]. Пошуковий запит один раз читає посилання і далі
працює з цим знімком до кінця, навіть якщо тим часом паралельна перебудова опублікувала
новий (схема build-then-swap). Ціною є те, що під час перебудови в пам'яті тимчасово
існують два індекси (підрозділ~3.10).

\subsection{Паралельна побудова індексу}

Побудова складається з п'яти фаз (рис.~\ref{fig:seq_index.png}):
\begin{enumerate}[label=\arabic*)]
\item \textit{list} (викликаючий потік) -- рекурсивний обхід каталогу, відбір \code{*.txt},
сортування шляхів; позиція у списку є ідентифікатором документа;
\item \textit{map} (пул) -- список ріжеться на батчі по 32 документи, тобто 525 задач для
16\,793 документів; задача читає файли і токенізує їх у приватну таблицю термів, нічого не
синхронізуючи;
\item \textit{merge} (пул, та сама задача) -- приватна таблиця передається будівнику
\code{IndexBuilder.addBatch}, який зливає її у спільну структуру;
\item \textit{finalize} (пул) -- по одній задачі на партицію сортує списки входжень і
перетворює їх на компактні масиви;
\item \textit{swap} -- новий знімок атомарно замінює попередній.
\end{enumerate}

Фіксований розмір батча зроблено свідомо: за будь-якої кількості потоків виконується однаковий
набір задач, тож у вимірюваннях змінюється лише паралелізм, а не гранулярність. 525 задач на
16 потоків -- понад 30 задач на потік, чого достатньо для балансування навантаження, коли
документи мають різну кількість унікальних слів. Довжини документів записуються в масив
\code{DocumentTable} без синхронізації: кожну комірку пише рівно одна задача, а
\code{Future.get()} робить записи видимими потоку, що виконує фінальну фазу. Неможливість
прочитати файл не перериває побудову: документ отримує нульову довжину і не враховується в
\code{documents\_indexed}; якщо не прочитано жодного документа, побудова відхиляється, а
чинний індекс зберігається.

\subsection{Неблокуюче злиття: партиції з CAS-скриньками}

\code{PartitionedIndexBuilder} -- основна реалізація спільної структури. Кожна з $P$ партицій
має приватну \code{HashMap} термів, неблокуючу «скриньку» (стек Трайбера [7]) і прапорець
власника \code{busy} типу \code{AtomicBoolean}. Задача, що завершила свій батч, розбиває його
на фрагменти за партиціями і для кожного фрагмента виконує такі кроки:
\begin{enumerate}[label=\arabic*)]
\item кладе фрагмент у скриньку партиції (операція CAS на вершині стека);
\item намагається стати власником партиції: \code{busy.compareAndSet(false, true)};
\item якщо вдалося -- вибирає зі скриньки всі фрагменти (свій і, можливо, чужі), зливає їх
у \code{HashMap} партиції, знімає прапорець і \textit{повторно перевіряє скриньку};
\item якщо партиція зайнята -- не чекає, а переходить до наступного фрагмента: її фрагмент
зіллє поточний власник (передача, hand-off).
\end{enumerate}

Така схема є різновидом «комбінування» (flat combining [8]): потік, що отримав доступ до
структури, виконує операції і за інших. Жоден потік ніколи не блокується на партиції.
Повторна перевірка після зняття прапорця закриває вікно, у якому фрагмент потрапив у скриньку
вже після останнього \code{pop} власника, але до \code{busy.set(false)}: той, хто поклав
фрагмент, бачить партицію зайнятою і йде далі, а власник після звільнення помічає непорожню
скриньку і знову бере її. Оскільки кожен, хто кладе фрагмент, після цього обов'язково
намагається стати власником, після повернення всіх викликів \code{addBatch} скриньки
порожні; фінальна фаза для певності виконує ще одне спорожнення.

Відсутність стану гонки обґрунтовується моделлю пам'яті Java [15]. \code{HashMap} партиції
змінює лише поточний власник. Право власності передається через volatile-прапорець:
\code{set(false)} має семантику release, успішний \code{compareAndSet} -- acquire, тому
кожен наступний власник бачить усі зміни попереднього. Фрагменти потрапляють від автора до
власника через CAS на вершині стека, що також встановлює happens-before. Проблема ABA для
стека Трайбера в Java не виникає: вузол не може бути повторно використаний, доки на нього
посилається хоч один потік, бо пам'ять звільняє збирач сміття. Взаємне блокування
неможливе, бо замків немає взагалі.

Важливо, що неблокуючий протокол прибирає \textit{очікування}, але не \textit{серіалізацію}:
злиття в одну партицію все одно виконується послідовно, тільки тим потоком, який першим
захопив власність. Наслідки цього видно з експерименту з однією партицією (підрозділ~3.7).

\subsection{Блокувальна альтернатива: lock striping}

Для порівняння реалізовано \code{StripedIndexBuilder} -- класичну схему розщеплення замків:
словник поділено на $S$ смуг, кожна з яких -- \code{HashMap} під власним \code{ReentrantLock}.
Задача групує батч за смугами і зливає кожну групу, утримуючи замок відповідної смуги. Потік
у кожен момент утримує щонайбільше один замок, тому кругове очікування, а з ним і взаємне
блокування, неможливі. За $S = 1$ схема вироджується в один глобальний замок. Будівник
вимірює кількість захоплень, кількість випадків, коли замок був зайнятий (спершу
\code{tryLock}, потім блокувальний \code{lock}), і сумарний час очікування -- ці лічильники
використано в підрозділі~3.7.

\subsection{Токенізація: дві версії}

Токен -- це максимальна послідовність ASCII-літер і цифр довжиною від 2 до 32 символів,
приведена до нижнього регістру; довші послідовності (хеші, злиті рядки) відкидаються, будь-який
інший байт, зокрема не-ASCII, є роздільником. Англійські службові слова (132 слова) у
індекс не потрапляють. Ті самі правила застосовуються до тексту запиту і до термів від моделі.

\textit{Перша версія} токенізатора створювала рядок \code{String} для кожного слова і
підраховувала частоти в \code{HashMap}, що очищалася для кожного документа. Для 1,6~ГБ
тексту це сотні мільйонів короткоживучих об'єктів. Розклад часу за фазами (підрозділ~3.5)
показав, що сумарний час токенізації зростає з кількістю потоків значно швидше, ніж мав би:
потоки конкурують за спільні ресурси -- виділення пам'яті, роботу збирача сміття, шину пам'яті.

\textit{Друга версія} (\code{BatchTermTable}) хешує байти слова на місці, не копіюючи їх, і
шукає терм у таблиці з відкритою адресацією та лінійним зондуванням, що належить задачі.
Хеш обчислюється тією самою формулою $h = 31h + c$, що й \code{String.hashCode()}, тож
порівняння з уже збереженим рядком виконується без перетворень. Рядок створюється лише тоді,
коли терм уперше трапляється в батчі. Частоти накопичуються в парі полів
«поточний документ / лічильник» запису і скидаються в список входжень при зміні документа,
тому окремої таблиці на документ не існує. Стару версію залишено як параметр
\code{"mapper": "hashmap"}; тести перевіряють, що обидві версії будують ідентичний індекс.
Ефект заміни наведено в підрозділі~3.6.

\subsection{HTTP-сервер і паралельна обробка запитів}

\code{HttpServer} написано поверх блокувальних сокетів. Окремий потік виконує цикл
\code{accept} і передає кожне з'єднання в пул HTTP-потоків. Потік пулу обслуговує з'єднання
протягом усього його життя: читає запит, викликає обробник, записує відповідь і, якщо клієнт
не просив закрити з'єднання, читає наступний запит (HTTP/1.1 keep-alive [14]). Повторне
використання TCP-з'єднання прибирає з кожного запиту рукостискання і є головною причиною
високої пропускної здатності в підрозділі~3.8. Тайм-аут простою з'єднання -- 15~с, після
10\,000 запитів сервер закриває з'єднання заголовком \code{Connection: close}, щоб один клієнт
не утримував потік безстроково. Запит на WebSocket-рукостискання переводить сокет у пул
WebSocket-сеансів, тож тривалі сеанси не займають HTTP-потоки.

Сервер обмежує вхідні дані: рядок запиту із заголовками -- до 16~КіБ (431), тіло -- до 1~МіБ
(413), тіло з \code{Transfer-Encoding} не підтримується (501). Невідомий шлях дає 404,
невідповідний метод -- 405 із заголовком \code{Allow}, звичайний GET на \code{/ws/search} --
426. Невідомі поля JSON ігноруються із записом у журнал. Під час зупинки (\code{docker stop}
надсилає SIGTERM, який обробляє shutdown hook JVM) сервер закриває слухаючий сокет і
простоюючі keep-alive з'єднання, а запитам, що вже виконуються, дає до 8~с на завершення.

\subsection{WebSocket}

Протокол RFC~6455 [13] реалізовано без сторонніх бібліотек. Рукостискання перевіряє метод
GET, заголовки \code{Connection: Upgrade}, \code{Sec-WebSocket-Key} (16 байт у Base64) і версію
13; відповідь містить \code{Sec-WebSocket-Accept = Base64(SHA-1(key + GUID))}, що перевірено на
прикладі з самого RFC. Читач кадрів підтримує довжини 7/16/64 біти, демаскування, ping/pong і
\textit{фрагментовані повідомлення} (кадри продовження). Порушення протоколу закривають
з'єднання з відповідним кодом: зарезервовані біти, незамаскований кадр, керуючий кадр понад
125 байтів, продовження без початого повідомлення -- 1002; бінарне повідомлення -- 1003;
некоректний UTF-8 -- 1007; повідомлення понад 64~КіБ -- 1009. Запис кадрів серіалізовано
замком, щоб pong, надісланий під час запису великого повідомлення, не вклинився в нього.
Закриття виконується за повним протоколом: сервер надсилає close-кадр і до 1~с чекає на
зустрічний.

\subsection{Ранжування та сніпети}

Документи ранжуються функцією Okapi BM25 [9]:
\begin{equation}
\mathrm{score}(d) = \sum_{t \in q} w_t \cdot \mathrm{idf}(t) \cdot
\frac{tf_{t,d}\,(k_1 + 1)}{tf_{t,d} + k_1\left(1 - b + b\,\dfrac{|d|}{\mathrm{avgdl}}\right)},
\end{equation}
\begin{equation}
\mathrm{idf}(t) = \ln\!\left(1 + \frac{N - df_t + 0{,}5}{df_t + 0{,}5}\right),
\end{equation}
де $k_1 = 1{,}2$, $b = 0{,}75$, $N$ -- кількість документів, $df_t$ -- документна частота
терма, $|d|$ -- довжина документа в термах. Вага $w_t$ дорівнює 1 для слів користувача і 0,4 для
слів, доданих моделлю: розширення підвищує повноту, але не може переважити те, що
користувач написав сам. На відміну від простого підрахунку збігів, BM25 враховує
рідкісність терма і частоту його вживання з поправкою на довжину документа, тож документ,
де «whale» трапляється сто разів, випереджає документ із випадковою згадкою.

Обчислення не потребує синхронізації: кожен запит виділяє власний масив-акумулятор на $N$
елементів, проходить списки входжень своїх термів і вибирає 10 найкращих документів
мін-купою. Рівні оцінки впорядковуються за ідентифікатором документа, тому видача
детермінована. Поле \code{score} у відповіді -- оцінка, поділена на найкращу оцінку запиту
(діапазон $(0; 1]$), сире значення BM25 повертається в полі \code{bm25}.

Сніпет будується читанням файла блоками по 64~КіБ до першого входження найрідкіснішого
знайденого слова користувача (якщо таких немає -- слова розширення). Збіг зараховується лише
за цілим словом, тому \code{parallel} не знаходиться всередині \code{unparalleled}, а сніпет
узгоджений з ранжуванням. Вікно до 180 символів зсувається до меж слів, пробіли згортаються,
не-ASCII байти замінюються символом-заповнювачем.

\subsection{Розширення запитів через LLM}

\code{QueryExpander} звертається до Ollama \code{/api/generate} без стрімінгу, з форматом
відповіді JSON, температурою 0,2 і лімітом 128 токенів. Промпт вимагає 5--8 англійських
синонімів або близьких ключових слів окремими словами, без повторів і без слів самого
запиту, і містить приклад «запит -- відповідь» на сторонньому запиті. Оскільки модель на CPU
є найповільнішою ланкою конвеєра, навколо виклику побудовано чотири механізми:
\begin{itemize}
\item \textit{LRU-кеш} останніх 1\,024 розширень (\code{LinkedHashMap} у порядку доступу під
замком);
\item \textit{об'єднання однакових запитів} (single-flight): одночасні запити з тим самим
набором термів чекають на один спільний \code{CompletableFuture}, зареєстрований атомарним
\code{putIfAbsent} у \code{ConcurrentHashMap};
\item \textit{обмеження паралельності}: не більше двох одночасних викликів моделі
(\code{Semaphore}, узгоджено з \code{OLLAMA\_NUM\_PARALLEL=2});
\item \textit{дедлайн}: на очікування слота, сам виклик і очікування спільного результату
разом дається 20~с; після тайм-ауту чи будь-якої помилки моделі пошук виконується за
початковим запитом, а у відповіді вказується джерело \code{timeout} або \code{error}.
\end{itemize}

Відповідь моделі розбирається захищено: прибираються обгортки коду, вирізається JSON-об'єкт
або масив з будь-якого оточення, а якщо відповідь обірвана лімітом токенів, терми
відновлюються регулярним виразом за закритими лапками. Отримані слова проходять той самий
токенізатор, що й індекс, з них прибираються дублікати, службові слова, числа і слова запиту;
залишається не більше 8 термів.

\subsection{Клієнти}

\textit{CLI-клієнт} (\code{SearchCli}) працює всередині контейнера, тому на хості не
потрібна Java: обгортки \code{./search} (Git Bash, Linux) і \code{search.cmd} (Windows)
викликають \code{docker compose exec server search}. Клієнт підтримує пошук з розширенням і
без, кількість результатів, побудову індексу та вивід статистики; у терміналі знайдені слова
підсвічуються кольором.

\textit{Web UI} -- одна сторінка на HTML, CSS і JavaScript без фреймворків, яку роздає сам
сервер. Пошук іде через WebSocket: три картки етапів конвеєра підсвічуються в міру надходження
повідомлень, показують час кожного етапу і джерело розширення, слова користувача та моделі
відображаються різними кольорами, у сніпетах вони виділяються. З'єднання перевідкривається
після розриву. Окрема панель будує індекс із заданою кількістю потоків і реалізацією та
показує статистику. Усі дані сервера вставляються в сторінку лише через \code{textContent} і
вузли DOM, без \code{innerHTML}, що виключає XSS через вміст документів.

На рис.~\ref{fig:webui.png} показано пошук «sherlock holmes» з розширенням: модель додала
терми «consulting», «detective», «baker», «street» (позначені синім), розширення тривало
734~мс, пошук за індексом -- 9~мс. Той самий запит без розширення (рис.~\ref{fig:webui_plain.png})
виконується за 2~мс; порядок результатів змінюється, бо документи, де трапляються також
слова розширення, отримують додаткову оцінку.

\fig{webui.png}{Web UI: пошук з розширенням через LLM}
\fig{webui_plain.png}{Web UI: той самий запит без розширення}
\IfFileExists{figures/cli.png}{\fig{cli.png}{CLI-клієнт: пошук з розширенням}}{}

\subsection{Контейнеризація}

\code{Dockerfile} багатоетапний: на етапі збирання в образі \code{eclipse-temurin:21-jdk}
скрипт \code{build.sh} компілює код, пакує jar і запускає всі тести -- якщо хоч один тест не
проходить, образ не збирається. Робочий образ на базі \code{eclipse-temurin:21-jre} містить
лише jar, статичні файли і скрипти. JVM запускається з
\code{-XX:MaxRAMPercentage=75 -XX:+UseG1GC} і стає процесом PID~1, тому отримує SIGTERM від
\code{docker stop}. Перевірка стану контейнера виконується без curl -- запитом через
\code{/dev/tcp} у bash.

\section{Тестування та результати}

\subsection{Показники масштабованості}

Нехай $T_1$ -- час виконання на одному потоці, $T_p$ -- на $p$ потоках. Прискорення та
ефективність визначаються як
\begin{equation}
S(p) = \frac{T_1}{T_p}, \qquad E(p) = \frac{S(p)}{p}.
\end{equation}

Закон Амдала [10] описує межу прискорення програми, частка $f$ якої виконується послідовно:
\begin{equation}\label{eq:amdahl}
S(p) = \frac{1}{f + (1 - f)/p}, \qquad \lim_{p \to \infty} S(p) = \frac{1}{f}.
\end{equation}
Параметр $f$ для експериментальних даних оцінено методом найменших квадратів за точками
$p \le 12$. Закон Густафсона [12] доповнює цю модель для випадку, коли разом із кількістю
процесорів зростає обсяг задачі; у роботі корпус фіксований, тож застосовується модель
Амдала.

Метрика Карпа -- Флатта [11] -- експериментально визначена послідовна частка:
\begin{equation}\label{eq:kf}
e(p) = \frac{1/S(p) - 1/p}{1 - 1/p}.
\end{equation}
Якщо $e$ не змінюється зі зростанням $p$, обмеження спричинене справді послідовною частиною
програми; якщо $e$ зростає -- накладними витратами паралелізму, що збільшуються разом із
кількістю потоків (синхронізація, конкуренція за спільні ресурси) [11].

\subsection{Методика та середовище}

Усі вимірювання виконано на машині автора в штатному розгортанні Docker Compose
(табл.~\ref{tab:env}). Вимірювання автоматизовано класом \code{Bench}, який запускається
всередині контейнера сервера (\code{bench/run\_all.sh}); сирі результати записуються у
CSV-файли каталогу \code{bench/results}, графіки будує \code{bench/plot.py}.

\begin{table}[H]
\caption{Середовище вимірювань}\label{tab:env}
\centering\begin{tabularx}{\textwidth}{lX}
\toprule
Компонент & Характеристика\\
\midrule
Процесор & AMD Ryzen 5 7500F: 6 ядер, 12 апаратних потоків (SMT), базова частота 3,7~ГГц,
L3 32~МБ\\
Пам'ять, доступна Docker & 15,7~ГіБ (віртуальна машина WSL2)\\
Платформа & Windows, Docker Desktop (WSL2), Docker Engine 28.3.3\\
JVM & Eclipse Temurin 21 (JRE), G1 GC, \code{-XX:MaxRAMPercentage=75}\\
LLM & Ollama, модель qwen2.5:1.5b, \code{OLLAMA\_NUM\_PARALLEL=2}, обчислення на CPU\\
Корпус & 16\,793 документи, 1,62~ГБ, том Docker на файловій системі ext4\\
\bottomrule
\end{tabularx}
\end{table}

Методика вимірювання часу побудови:
\begin{itemize}
\item перед серією виконується прогрівна побудова (JIT-компіляція, кеш сторінок ОС), її
результат не враховується;
\item кожна конфігурація вимірюється тричі; три прогони йдуть раундами, і в кожному раунді
порядок конфігурацій випадково перемішується (фіксоване зерно), щоб повільні дрейфи -- нагрів
процесора, стан кешу -- розподілялися рівномірно, а не зміщували окрему конфігурацію;
\item після кожної побудови викликається \code{System.gc()}, щоб сміття попередньої побудови
не збиралося під час наступної;
\item час вимірюється на сервері від початку обходу каталогу до завершення фінальної фази;
у таблицях наведено медіану трьох прогонів, на графіках -- також мінімум і максимум.
\end{itemize}

Пропускну здатність вимірює вбудований навантажувальний клієнт із замкненим циклом: кожен із
$c$ потоків-клієнтів тримає власне keep-alive з'єднання і надсилає наступний запит одразу
після отримання відповіді. Запити циклічно беруться з набору 20 запитів різної частотності
(\code{bench/queries.txt}). Після 5~с прогріву рахуються запити протягом 20~с (для LLM --
30~с), латентність фіксується для кожного запиту окремо. Обмеження методики: клієнт працює в
тому самому контейнері, що й сервер, і ділить з ним 12 апаратних потоків, тому виміряна
пропускна здатність є нижньою оцінкою можливостей сервера.

\subsection{Модульне, стресове та наскрізне тестування}

Для тестів написано власний мінімальний раннер (анотація \code{@Test}, перевірки
\code{check}/\code{equal}/\code{throwsType}); тести запускаються командою
\code{./build.sh test} і під час кожного \code{docker build}. Усього 38 тестів у 9 класах
(табл.~\ref{tab:tests}), повний прогін триває близько 10~с.

\begin{longtable}{p{0.27\textwidth}p{0.66\textwidth}}
\caption{Склад тестів}\label{tab:tests}\\
\toprule
Клас & Що перевіряється\\
\midrule\endfirsthead
\multicolumn{2}{l}{\textit{Продовження таблиці~\ref{tab:tests}}}\\
\toprule Клас & Що перевіряється\\ \midrule\endhead
\code{ThreadPoolTest} & результати через \code{Future}; справжня паралельність (4 задачі
чекають одна на одну на \code{CountDownLatch}); передача винятків і живучість потоків;
graceful shutdown виконує всі 1\,000 поставлених задач і відхиляє нові; \code{shutdownNow}
повертає невиконані; 8 потоків-продюсерів ставлять 80\,000 задач -- кожна виконується рівно
один раз\\
\code{TreiberStackTest} & 8 потоків по 50\,000 операцій push/pop: жоден елемент не втрачено і не
вийнято двічі\\
\code{IndexBuilderTest} & індекс, побудований паралельно, збігається з послідовним еталоном
(окремий простий алгоритм) для обох токенізаторів, обох реалізацій, 1--16 потоків і
1/3/64 партицій -- 60 побудов; 15 побудов у режимі найбільшої конкуренції (батч з одного
документа, 1--2 партиції, 12 потоків); незалежність ідентифікаторів документів від кількості
потоків; колізії хешів (\code{"Aa"} і \code{"BB"}), регістр, службові слова; порожній корпус\\
\code{TokenizerTest} & регістр, роздільники, довжини, не-ASCII, кирилиця, дедуплікація\\
\code{SearchTest} & BM25: вища частота і рідкісніший терм ранжуються вище, слова моделі
важать менше, детермінізм за рівних оцінок; сніпет за цілим словом, на межі блоків читання\\
\code{QueryExpanderTest} & розбір коректних, «загорнутих» і обірваних відповідей моделі;
очищення термів; кеш; деградація за недоступної моделі; дедлайн 300~мс за відповіді
моделі через 3~с; 12 одночасних запитів (4 різні) -- не більше 4 викликів моделі і не більше
2 одночасно (mock-сервер Ollama)\\
\code{JsonTest} & вкладені значення, екранування, некоректні входи\\
\code{WebSocketTest} & ключ рукостискання за прикладом RFC~6455\\
\code{IntegrationTest} & наскрізні тести на запущеному сервері: контракт усіх ендпоінтів і коди
помилок; 50 запитів в одному keep-alive з'єднанні; 16 клієнтів по 20 пошуків
\textit{під час перебудови індексу} -- усі 320 відповідей ідентичні; WebSocket: три етапи,
два запити в одному сеансі, фрагментоване повідомлення, ping/pong, етап \code{error}, коди
закриття 1000/1002/1003/1007; пошук за повільної моделі завершується в межах дедлайну\\
\bottomrule
\end{longtable}

У Java немає аналога ThreadSanitizer, тому відсутність станів гонки перевіряється стресовими
тестами з порівнянням з еталоном: будь-яка втрата чи дублювання входжень під час
конкурентного злиття змінила б частоти термів і була б виявлена. Тест із перебудовою під
навантаженням перевіряє коректність публікації знімка: жоден пошук не побачив частково
побудованого індексу.

\subsection{Час побудови індексу залежно від кількості потоків}

Результати для обох реалізацій спільної структури з 64 партиціями наведено в
табл.~\ref{tab:index} і на рис.~\ref{fig:index_scaling.png}, \ref{fig:index_efficiency.png}.

\begin{table}[H]
\caption{Час побудови індексу (64 партиції, медіана трьох прогонів)}\label{tab:index}
\centering\small
\begin{tabular}{c|ccccc|cc}
\toprule
& \multicolumn{5}{c|}{partitioned (неблокуюча)} & \multicolumn{2}{c}{striped (блокувальна)}\\
$p$ & $T_p$, с & min--max, с & $S$ & $E$ & $e$ & $T_p$, с & $S$\\
\midrule
\tabinput{tables/index.tex}
\bottomrule
\end{tabular}
\end{table}

\fig{index_scaling.png}{Час побудови індексу та прискорення залежно від кількості потоків}
\fig{index_efficiency.png}{Ефективність та експериментальна послідовна частка Карпа -- Флатта}

На двох -- шести потоках прискорення близьке до лінійного з невеликими втратами: 1,75 на 2
потоках і 4,02 на 6 (ефективність 0,87 і 0,67). Шість -- кількість фізичних ядер процесора; за
більшої кількості потоків два потоки ділять одне ядро (SMT), і приріст різко сповільнюється:
від 6 до 10 потоків прискорення зростає лише з 4,02 до 4,85, далі не зростає зовсім, а на
16 потоках, коли потоків більше, ніж апаратних потоків процесора, -- трохи падає (4,51) через
витрати на перемикання контексту. Найкращий результат -- 2,09~с на 10 потоках.

Апроксимація законом Амдала дає послідовну частку $f = 0{,}125$, тобто межу прискорення
$1/f = 8$. Проте справді послідовні ділянки програми значно менші: обхід каталогу триває
близько 78~мс, а фінальна фаза -- близько 0,1~с (табл.~\ref{tab:phases}), разом менше 2~\%
від $T_1$. Експериментальна частка Карпа -- Флатта не постійна: вона становить 0,08--0,14 до
12 потоків і зростає до 0,17 на 16. За критерієм підрозділу~3.1 це означає, що масштабування
обмежують накладні витрати, які зростають разом із кількістю потоків, а не послідовний код.
Їхню природу показує розклад за фазами.

Різниця між реалізаціями спільної структури за 64 партицій лежить у межах розкиду прогонів:
на 8 і 16 потоках трохи швидша блокувальна, на 6, 10 і 12 -- неблокуюча. За такої кількості
партицій синхронізація не є вузьким місцем (підрозділ~3.7).

\subsection{Розклад часу за фазами}

\begin{table}[H]
\caption{Медіани фаз побудови (partitioned, 64 партиції)}\label{tab:phases}
\centering\small
\begin{tabular}{cccccccc}
\toprule
& \multicolumn{3}{c}{Стіновий час фаз} & \multicolumn{3}{c}{Сумарний час потоків, с} &
Передач\\
$p$ & list, мс & map+merge, с & finalize, мс & читання & токенізація & злиття & фрагментів\\
\midrule
\tabinput{tables/phases.tex}
\bottomrule
\end{tabular}
\end{table}

\fig[0.85\textwidth]{index_phases_cpu.png}{Сумарний час потоків за фазами задачі}

Сумарний (за всіма потоками) час кожної фази задачі показано в табл.~\ref{tab:phases} і на
рис.~\ref{fig:index_phases_cpu.png}. Якби накладних витрат не було, висота стовпчиків
залишалася б сталою: та сама робота лише розподіляється між більшою кількістю потоків.
Насправді сумарний час токенізації зростає з 6,84~с на одному потоці до 9,49~с на шести і
15,41~с на дванадцяти, тобто в 2,25 раза. Токенізація не має жодної синхронізації, тож це
зростання -- ознака конкуренції за апаратні ресурси: за SMT два потоки ділять обчислювальні
блоки й кеші L1/L2 одного ядра, усі потоки -- спільний кеш L3 і шину пам'яті, а частота
процесора за навантаження всіх ядер нижча, ніж за одного. Злиття зростає з 2,47 до 5,99~с з
тих самих причин плюс через частіші промахи кешу під час вставок у спільні хеш-таблиці.
Читання файлів займає 0,5--1,8~с сумарно і вузьким місцем не є: після перенесення корпусу в том
Docker (підрозділ~2.1) задача стала обчислювальною.

\subsection{Оптимізація токенізації}

Перша версія задачі створювала об'єкт \code{String} і робила пошук у \code{HashMap} для
кожного слова. Розклад за фазами для неї показав ще стрімкіше зростання сумарного часу
токенізації -- з 10,78~с до 26,13~с на 12 потоках (×2,42) і 37,40~с на 16 потоках. Друга
версія (\code{BatchTermTable}, підрозділ~2.7) хешує байти на місці й створює рядок лише для
нових у батчі термів. Порівняння виконано в тій самій сесії тим самим інструментом
(табл.~\ref{tab:mapper}, рис.~\ref{fig:mapper_compare.png}).

\begin{table}[H]
\caption{Порівняння версій токенізації (partitioned, 64 партиції)}\label{tab:mapper}
\centering\begin{tabular}{cccccc}
\toprule
& \multicolumn{2}{c}{v1: String на слово} & \multicolumn{2}{c}{v2: байтова таблиця} & \\
$p$ & $T_p$, с & токенізація, с & $T_p$, с & токенізація, с & $T^{v1}_p / T^{v2}_p$\\
\midrule
\tabinput{tables/mapper.tex}
\bottomrule
\end{tabular}
\end{table}

\fig{mapper_compare.png}{Час побудови і прискорення двох версій токенізації}

Нова версія швидша в 1,28--1,48 раза за будь-якої кількості потоків: $T_1$ скоротився з
14,28 до 10,12~с, а найкращий час -- з 2,95 до 2,09~с. Сумарний час токенізації зменшився
приблизно вдвічі (на 12 потоках -- з 26,13 до 15,41~с). Прискорення відносно власного $T_1$ в
обох версій майже однакове (рис.~\ref{fig:mapper_compare.png}, праворуч), тобто оптимізація
зменшила обсяг роботи на кожному потоці, але не усунула апаратну конкуренцію, описану в
підрозділі~3.5. Цей результат ілюструє загальне правило: зменшення обсягу виділення пам'яті в
гарячому циклі дає виграш, який множиться на всі потоки, тоді як заміна механізму
синхронізації за малої конкуренції не дає майже нічого (підрозділ~3.7).

\subsection{Вплив кількості партицій і механізму синхронізації}

Щоб навмисно створити конкуренцію за спільну структуру, побудову повторено з 1, 8 і 256
партиціями (смугами) для обох реалізацій (табл.~\ref{tab:partitions},
рис.~\ref{fig:index_partitions.png}, \ref{fig:lock_wait.png}).

\begin{table}[H]
\caption{Час побудови (с) залежно від кількості партицій; синхронізація на 12 потоках}
\label{tab:partitions}
\centering\small
\begin{tabular}{llcccccl}
\toprule
& & \multicolumn{4}{c}{$T_p$ за $p$ потоків, с} & Злиття & Синхронізація\\
Реалізація & $P$ & 1 & 4 & 8 & 12 & (12), с & (12 потоків)\\
\midrule
\tabinput{tables/partitions.tex}
\bottomrule
\end{tabular}
\end{table}

\fig{index_partitions.png}{Час побудови за різної кількості партицій}
\fig[0.8\textwidth]{lock_wait.png}{Сумарний час очікування блокувань (striped)}

З однією партицією обидві реалізації перестають масштабуватися вже після 4 потоків:
3,4--3,7~с на 8 і 12 потоках проти 2,1~с з 8 і більше партиціями. Але причини різні.
Блокувальна реалізація з одним замком на 12 потоках сумарно чекає на нього 22,4~с -- потоки
простоюють, а сумарний час фази злиття зростає до 26,1~с. Неблокуюча реалізація не чекає
зовсім: зайнятість партиції фіксувалася 477 разів, і щоразу фрагмент передавався власнику, а
потік ішов далі. Проте вся робота злиття все одно виконується послідовно -- тим потоком, що
володіє єдиною партицією, -- тож стіновий час той самий. Неблокуючий алгоритм гарантує
прогрес і відсутність простою, але не робить послідовну за природою операцію паралельною: це
наочна ілюстрація закону Амдала, де послідовною частиною є злиття в одну структуру.

Вже 8 партицій прибирають цей ефект для обох реалізацій: очікування на замках падає до
1,5~с (на 12 потоків за 2,1~с роботи), час побудови -- до 2,1~с. Подальше збільшення
кількості партицій до 64 і 256 не прискорює побудову, а в неблокуючій реалізації збільшує
кількість фрагментів і передач (до 6\,330 на 256 партицій), що лише трохи додає накладних
витрат. Отже, за достатньої кількості партицій механізм синхронізації спільної структури
перестає впливати на результат, і масштабування визначається обчислювальною частиною задачі
(підрозділ~3.5). За замовчуванням використовується неблокуюча реалізація з 64 партиціями:
вона ніколи не блокує потік, а її продуктивність не гірша за блокувальну.

\subsection{Пропускна здатність за паралельних клієнтів}

\begin{table}[H]
\caption{Пропускна здатність пошуку без розширення (20~с після 5~с прогріву)}
\label{tab:throughput}
\centering\begin{tabular}{cccccc}
\toprule
Клієнтів & Запитів/с & p50, мс & p95, мс & p99, мс & Помилок\\
\midrule
\tabinput{tables/throughput.tex}
\bottomrule
\end{tabular}
\end{table}

\fig{throughput.png}{Пропускна здатність і латентність пошуку залежно від кількості клієнтів}

Пропускна здатність (табл.~\ref{tab:throughput}, рис.~\ref{fig:throughput.png}) зростає
майже лінійно до 8 клієнтів: 1\,247 запитів/с на одному, 9\,172 на восьми (×7,4) за майже
незмінної медіанної латентності 0,76--0,81~мс. Це очікувано: пошукові потоки лише читають
незмінний знімок індексу, не мають спільних записуваних даних і не синхронізуються між собою.
Максимум -- 12\,004 запитів/с -- досягається на 16 клієнтах; далі пропускна здатність
стабілізується на рівні 11\,500, а латентність зростає (p99 -- 13,1~мс на 32 і 31,7~мс на 64
клієнтах): усі 12 апаратних потоків зайняті, і нові клієнти лише довше чекають на процесор.
Частину процесорного часу при цьому споживає сам навантажувальний клієнт, що працює на тій
самій машині. Помилок не зафіксовано.

\subsection{Вузьке місце LLM-розширення}

\begin{table}[H]
\caption{Пошук з розширенням через LLM}\label{tab:llm}
\centering\begin{tabular}{lcccc}
\toprule
Режим & Клієнтів & Запитів/с & p50, мс & p95, мс\\
\midrule
\tabinput{tables/llm.tex}
\bottomrule
\end{tabular}
\end{table}

\fig{throughput_llm.png}{Пропускна здатність і латентність пошуку з розширенням (без кешу)}

Виклик моделі на CPU триває близько 0,6~с, що на три порядки довше за сам пошук
(табл.~\ref{tab:llm}, рис.~\ref{fig:throughput_llm.png}). Без кешу пропускна здатність
становить 1,6~запиту/с для одного клієнта і 2,3--2,4 запиту/с для двох і більше: семафор
пропускає до моделі не більше двох запитів одночасно, а Ollama обробляє не більше двох
паралельно. Подальше збільшення кількості клієнтів лише подовжує чергу, і латентність зростає
пропорційно: 0,89~с для двох, 1,69~с для чотирьох, 3,32~с для восьми клієнтів. Це узгоджується
із законом Літтла $L = \lambda W$: за $\lambda \approx 2{,}4$~запиту/с і $L = 8$ запитів у
системі $W \approx 3{,}3$~с. Модель -- типовий послідовний ресурс конвеєра, і вирішальним є
не те, як паралельно працює решта системи, а те, щоб очікування на модель не займало
спільні ресурси. Тому в реалізації (підрозділ~2.11) очікування обмежене дедлайном, знімок
індексу береться лише після розширення, а однакові запити об'єднуються.

Кешування розширень повністю прибирає модель із шляху повторних запитів: з кешем пропускна
здатність становить 959 запитів/с на одному клієнті і 8\,634 на 32 клієнтах, що близько до
пошуку без розширення (решту часу займають пошук за розширеним набором термів і кеш під
замком).

\subsection{Поведінка пам'яті}

\begin{table}[H]
\caption{Пам'ять під час шести послідовних перебудов на 12 потоках, МіБ}\label{tab:memory}
\centering\begin{tabular}{c|ccc|ccc}
\toprule
& \multicolumn{3}{c|}{без \code{System.gc()}} & \multicolumn{3}{c}{з \code{System.gc()}}\\
Перебудова & heap used & committed & RSS & heap used & committed & RSS\\
\midrule
\tabinput{tables/memory.tex}
\bottomrule
\end{tabular}
\end{table}

\fig{memory.png}{Пам'ять JVM під час повторних перебудов індексу}

Готовий індекс займає в купі близько 350~МіБ (табл.~\ref{tab:memory}, рис.~\ref{fig:memory.png}):
33,2~млн упакованих входжень по 8~байт (253~МіБ) плюс словник. Без явного збирання сміття
після перебудови зайнята купа становить 1,7--3,2~ГіБ: G1 не поспішає збирати старий індекс і
тимчасові структури побудови, доки є вільна пам'ять. З викликом \code{System.gc()} після підміни
зайнята купа стабільно повертається до 352~МіБ, тобто витоку пам'яті між перебудовами немає.
RSS процесу при цьому коливається від 2,5 до 5,2~ГіБ: JVM утримує раз виділені сторінки купи
(\code{committed} -- 1,3~ГіБ після збирання, але пікові сторінки під час побудови повертаються
операційній системі лише частково). Звідси практичний висновок для розгортання: на час
перебудови потрібно передбачити пам'ять щонайменше для двох індексів і проміжних таблиць
задач, тобто кілька гігабайтів, хоча сам індекс займає лише 350~МіБ.

\subsection{Функціональні тестові запити}

Для перевірки якості пошуку виконано набір запитів, очікуваний результат яких відомий
наперед (табл.~\ref{tab:queries}; ідентифікатор документа -- номер книжки Project Gutenberg і
номер частини). Без розширення пошук займає менше 1~мс, з розширенням -- 0,55--0,82~с.

\begin{longtable}{p{0.18\textwidth}p{0.25\textwidth}p{0.30\textwidth}p{0.14\textwidth}}
\caption{Результати пошуку на тестових запитах}\label{tab:queries}\\
\toprule
Запит & Очікуваний результат & Фактичний результат (топ-3) & Оцінка\\
\midrule\endfirsthead
\multicolumn{4}{l}{\textit{Продовження таблиці~\ref{tab:queries}}}\\
\toprule Запит & Очікуваний & Фактичний (топ-3) & Оцінка\\ \midrule\endhead
moby dick white whale & книжки «Moby Dick» (PG 15, 2489, 2701) & 15\_0000, 2489\_0003,
2701\_0003 -- три видання «Moby Dick» & відповідає\\
sherlock holmes baker street & оповідання про Шерлока Холмса & 2038\_0001 (збірка з «A Scandal
in Bohemia»), 221\_0000 і 108\_0000 («The Return of Sherlock Holmes») & відповідає\\
tarzan apes jungle & книжки про Тарзана (PG 78, 106) & 78\_0001 «Tarzan of the Apes»,
106\_0002, 106\_0003 «Jungle Tales of Tarzan» & відповідає\\
romeo juliet & п'єса «Romeo and Juliet» (PG 1513) & 1513\_0000, 1513\_0001, 1286\_0004 &
відповідає\\
pride prejudice darcy & «Pride and Prejudice» (PG 1342) & 1342\_0006, 1342\_0003, 1342\_0000 &
відповідає\\
abraham lincoln & біографії Лінкольна & 1815\_0000 «The Boys' Life of Abraham Lincoln»,
1713\_0006, 2517\_0006 & відповідає\\
electricity magnetism & праці з електромагнетизму & 1225\_0001, 1225\_0000 «Faraday As A
Discoverer», 1707\_0004 & відповідає\\
treasure island & «Treasure Island» (PG 120) & 973\_0000 «Howard Pyle's Book of Pirates»,
3029\_0000, 2021\_0009 & PG 120 немає в корпусі; знайдено тематично близьке\\
quixotic & рідкісне слово, точні збіги & 2158\_0012 «The Prime Minister», 1803\_0003,
2158\_0011 & відповідає\\
xyzzyplugh & порожня видача & без розширення -- порожня; з розширенням модель вигадала
«xylophone», «clavier», «melody» і дала нерелевантні документи & див. текст\\
the and of & порожня видача (лише службові слова) & порожня, звернення до моделі не
виконується & відповідає\\
\bottomrule
\end{longtable}

З розширенням видача змінюється змістовно: для «moby dick white whale» модель додала
«captain», «ahab», для «abraham lincoln» -- «president», «union», «republican», і на першому
місці опинився том праць самого Лінкольна (PG 2653); для «treasure island» -- «pirate»,
«captain», «gold», і у видачу потрапив «Peter Pan» (PG 16). Слабке місце розширення
показує безглуздий запит «xyzzyplugh»: модель не відмовляється від відповіді, а вигадує
асоціації за схожістю написання, і запит без жодного збігу отримує непорожню видачу.
Зменшена вага слів моделі (0,4) не допомагає, коли слів користувача в корпусі немає зовсім;
можливе покращення -- не розширювати запит, якщо жоден його терм не знайдено в індексі.

\unsection{ВИСНОВКИ}

У курсовій роботі розроблено контейнеризовану пошукову систему, що демонструє паралелізм на
всіх рівнях: паралельну побудову інвертованого індексу власним пулом потоків, потокобезпечну
спільну структуру індексу у двох варіантах -- неблокуючому та блокувальному, -- безблокувальне
читання незмінного знімка індексу і паралельну обробку запитів багатьох клієнтів власним
HTTP/1.1- і WebSocket-сервером. Пошуковий конвеєр розширює запит через локальну модель
qwen2.5:1.5b і ранжує документи функцією BM25. Система запускається командою
\code{docker compose up}, повністю реалізує фіксований контракт REST API і WebSocket, має два
клієнти (CLI і Web UI) і написана мовою Java 21 без жодних сторонніх бібліотек. Її
коректність підтверджують 38 тестів, зокрема стресові порівняння паралельних побудов з
послідовним еталоном і наскрізні тести під навантаженням.

Результати дослідження масштабованості на процесорі з 6 ядрами і 12 апаратними потоками:
\begin{enumerate}[label=\arabic*)]
\item Час побудови індексу корпусу з 16\,793 документів (1,62~ГБ) скорочується з 10,12~с на
одному потоці до 2,09~с на десяти: прискорення 4,85. До шести потоків (кількість фізичних
ядер) ефективність становить 0,67--0,87, після шести додаткові потоки використовують SMT і
дають лише 20~\% приросту, а на 16 потоках час трохи зростає.
\item Апроксимація законом Амдала дає послідовну частку 0,125, тоді як справді послідовні
ділянки займають менше 2~\% часу; експериментальна частка Карпа -- Флатта зростає з
кількістю потоків. Отже, масштабування обмежують накладні витрати, що зростають із
паралелізмом: сумарний час токенізації, яка не має жодної синхронізації, на 12 потоках у
2,25 раза більший, ніж на одному. Причина -- конкуренція за спільні апаратні ресурси (ядра
при SMT, кеш L3, шина пам'яті, частота процесора).
\item Заміна токенізації, що створювала рядок на кожне слово, байтовою хеш-таблицею задачі
прискорила побудову в 1,28--1,48 раза за будь-якої кількості потоків. Зменшення виділення
пам'яті в гарячому циклі дало більший ефект, ніж будь-яка зміна механізму синхронізації.
\item З однією партицією обидві реалізації спільної структури перестають масштабуватися після
4 потоків (3,4--3,7~с проти 2,1~с). Блокувальна при цьому сумарно простоює на замку 22,4~с,
неблокуюча не простоює зовсім, але виконує злиття тим самим послідовним потоком-власником:
неблокуючий алгоритм усуває очікування, а не серіалізацію. Від 8 партицій синхронізація
перестає впливати на час, і обидві реалізації показують однаковий результат.
\item Розміщення корпусу в томі Docker замість каталогу диска Windows скоротило побудову на
12 потоках з 42,3 до 4,0~с: через міст файлових систем Windows -- WSL2 задача була обмежена
введенням-виведенням і не могла показати масштабування.
\item Пропускна здатність пошуку зростає майже лінійно до 8 клієнтів і досягає 12\,004
запитів/с на 16 клієнтах за медіанної латентності менше 1~мс; далі вона обмежена кількістю
апаратних потоків, а латентність зростає. Пошукові потоки не синхронізуються між собою, бо
читають незмінний знімок індексу.
\item Розширення запиту через LLM на CPU обмежує пропускну здатність рівнем 2,3--2,4 запиту/с
незалежно від кількості клієнтів: модель є послідовним ресурсом, і латентність зростає за
законом Літтла. Кеш розширень повертає пропускну здатність повторних запитів до
8\,634 запитів/с, а дедлайн і підміна знімка після розширення не дають очікуванню на
модель утримувати ресурси сервера.
\item Готовий індекс займає близько 350~МіБ купи, витоку пам'яті між перебудовами немає, але
під час перебудови за схемою build-then-swap потрібна пам'ять для двох індексів і проміжних
таблиць, тож пікове використання становить кілька гігабайтів.
\end{enumerate}

Подальші напрями покращення: зменшити обсяг пам'яті, що виділяється під час побудови
(повторно використовувати таблиці задач, стискати списки входжень дельта-кодуванням), що
зменшить конкуренцію за шину пам'яті; не розширювати запит, якщо жоден його терм не знайдено
в індексі; винести модель на GPU або обробляти розширення пакетами, щоб послабити найбільше
послідовне обмеження конвеєра; врахувати в ранжуванні близькість слів запиту в документі.

\unsection{ПЕРЕЛІК ДЖЕРЕЛ ПОСИЛАННЯ}

\begin{enumerate}[label=\arabic*.,leftmargin=*,labelindent=\parindent,itemsep=1pt]
\makeatletter
\let\olditem\item
\renewcommand{\item}{\stepcounter{refs}\olditem}
\makeatother
\item Семеренко В. П. Технології паралельних обчислень : навч. посіб. Вінниця : ВНТУ, 2018.
105~с.
\item Мартін Р. Чистий код: створення, аналіз і рефакторинг. Харків : Фабула, 2019. 416~с.
\item Manning C. D., Raghavan P., Schütze H. Introduction to Information Retrieval. Cambridge :
Cambridge University Press, 2008. 506~p.
\item Goetz B., Peierls T., Bloch J. та ін. Java Concurrency in Practice. Boston :
Addison-Wesley, 2006. 432~p.
\item Williams A. C++ Concurrency in Action: Practical Multithreading. Shelter Island :
Manning, 2012. 528~p.
\item Herlihy M., Shavit N., Luchangco V., Spear M. The Art of Multiprocessor Programming.
2nd ed. Cambridge, MA : Morgan Kaufmann, 2020. 576~p.
\item Treiber R. K. Systems Programming: Coping with Parallelism : Research Report RJ~5118.
San Jose : IBM Almaden Research Center, 1986. 46~p.
\item Hendler D., Incze I., Shavit N., Tzafrir M. Flat Combining and the
Synchronization-Parallelism Tradeoff // Proc. 22nd ACM Symposium on Parallelism in
Algorithms and Architectures (SPAA). 2010. P.~355--364.
\item Robertson S., Zaragoza H. The Probabilistic Relevance Framework: BM25 and Beyond //
Foundations and Trends in Information Retrieval. 2009. Vol.~3, No.~4. P.~333--389.
\item Amdahl G. M. Validity of the Single Processor Approach to Achieving Large Scale
Computing Capabilities // AFIPS Spring Joint Computer Conference. 1967. P.~483--485.
\item Karp A. H., Flatt H. P. Measuring Parallel Processor Performance // Communications of
the ACM. 1990. Vol.~33, No.~5. P.~539--543.
\item Gustafson J. L. Reevaluating Amdahl's Law // Communications of the ACM. 1988. Vol.~31,
No.~5. P.~532--533.
\item Fette I., Melnikov A. The WebSocket Protocol : RFC~6455. IETF, 2011. URL:
\url{https://www.rfc-editor.org/rfc/rfc6455}.
\item Fielding R., Nottingham M., Reschke J. HTTP/1.1 : RFC~9112. IETF, 2022. URL:
\url{https://www.rfc-editor.org/rfc/rfc9112}.
\item Gosling J., Joy B., Steele G. та ін. The Java Language Specification, Java SE 21
Edition. Chapter 17: Threads and Locks. Oracle, 2023. URL:
\url{https://docs.oracle.com/javase/specs/jls/se21/html/jls-17.html}.
\item Ollama: API Reference [Електронний ресурс]. URL:
\url{https://github.com/ollama/ollama/blob/main/docs/api.md}.
\item Docker Documentation: Docker Compose [Електронний ресурс]. URL:
\url{https://docs.docker.com/compose/}.
\item Conventional Commits 1.0.0 [Електронний ресурс]. URL:
\url{https://www.conventionalcommits.org/}.
\item Google Java Style Guide [Електронний ресурс]. URL:
\url{https://google.github.io/styleguide/javaguide.html}.
\item Project Gutenberg [Електронний ресурс]. URL: \url{https://www.gutenberg.org/}.
\item Qwen Team. Qwen2.5 Technical Report : arXiv:2412.15115. 2024. URL:
\url{https://arxiv.org/abs/2412.15115}.
\end{enumerate}

\clearpage
\renewcommand{\thefigure}{А.\arabic{figure}}
\setcounter{figure}{0}
\section*{ДОДАТОК А}
\begin{center}\bfseries Детальні діаграми класів\end{center}
\addcontentsline{toc}{section}{ДОДАТОК А Детальні діаграми класів}

\fig[\textwidth]{class_core.png}{Діаграма класів пакета \code{core}}
\fig[\textwidth]{class_net_api.png}{Діаграма класів пакетів \code{net}, \code{api},
\code{llm}, \code{search}}

\clearpage
\section*{ДОДАТОК Б}
\begin{center}\bfseries Репозиторій і відтворення результатів\end{center}
\addcontentsline{toc}{section}{ДОДАТОК Б Репозиторій і відтворення результатів}

Вихідний код, історія комітів, UML-джерела, сирі результати вимірювань і ця записка
зберігаються в репозиторії:
\url{https://github.com/d1aboli/course_work_parallel_computing}.

Відтворення всіх результатів (Git Bash, у корені репозиторію, архів \code{D1.7GB.zip} поруч):
\begin{lstlisting}
docker compose up --build -d        # dataset-init, ollama-init, ollama, server
./search --index 12                 # побудова індексу
./search --expand "storm at sea"    # перевірка пошуку
bash bench/run_all.sh               # дослідження, результати в bench/results
python bench/plot.py                # графіки в report/figures
python report/gen_tables.py         # таблиці записки
cd report && xelatex main.tex && xelatex main.tex
\end{lstlisting}

Структура репозиторію:
\begin{lstlisting}
src/main/java/ua/kpi/pocw/
  core/    ThreadPool, Indexer, IndexBuilder, PartitionedIndexBuilder,
           StripedIndexBuilder, TreiberStack, InvertedIndex, BatchTermTable
  net/     HttpServer, HttpRequest, HttpResponse, WebSocket
  api/     Router, SearchService, WsSearchHandler
  search/  Searcher (BM25), Snippets
  llm/     OllamaClient, QueryExpander
  cli/     SearchCli           bench/  Bench
src/test/java/...              власний тест-раннер і 38 тестів
web/                           Web UI
docs/uml/                      PlantUML-діаграми
tools/prepare_dataset.py       підготовка корпусу
bench/                         сценарій дослідження, результати, графіки
report/                        пояснювальна записка (LaTeX)
\end{lstlisting}

\clearpage
\section*{ДОДАТОК В}
\begin{center}\bfseries Промпти, використані під час виконання роботи\end{center}
\addcontentsline{toc}{section}{ДОДАТОК В Промпти, використані під час виконання роботи}

\textit{Промпт 1. Розробка системи}

Мені потрібно зробити курсову з паралельних обчислень: пошукову систему, яка будує
інвертований індекс з текстових файлів у кілька потоків і шукає по ньому. Писати хочу на
Java, без сторонніх бібліотек: пул потоків і структуру індексу для одночасного запису з
багатьох потоків треба написати самому. Система має запускатися однією командою
\code{docker compose up} разом з Ollama і мати REST API та WebSocket за специфікацією з
силабусу, який я додаю.

Поясни, як влаштувати індекс, щоб потоки не заважали один одному, і зроби два варіанти: з
блокуваннями і без них, щоб їх можна було порівняти. Додай простий веб-інтерфейс і консольний
клієнт, тести на відсутність станів гонки й README з інструкцією запуску.

\bigskip
\textit{Промпт 2. Дослідження швидкодії}

Система вже працює на моєму комп'ютері (Ryzen 5 7500F, 6 ядер, 12 потоків). Допоможи
виміряти, як час побудови індексу залежить від кількості потоків від 1 до 16. Кожен варіант
запусти по три рази в перемішаному порядку, щоб результати були чесними. Окремо потрібно
виміряти, скільки запитів за секунду витримує пошук при різній кількості одночасних клієнтів, з
розширенням запиту через LLM і без нього.

Побудуй з результатів графіки й поясни простими словами, де вузьке місце і чому прискорення не
росте далі. Якщо по вимірах видно, що щось можна оптимізувати, зроби це й порівняй «до» і
«після».


\end{document}
